% !TeX program = lualatex
\documentclass[11pt,openany]{book}

\usepackage{fontspec}
\usepackage[brazil]{babel}
\usepackage{microtype}
\usepackage[a4paper,margin=2.6cm,headheight=15pt]{geometry}
\usepackage{graphicx}
\usepackage{booktabs}
\usepackage{longtable}
\usepackage{amsmath}
\usepackage[dvipsnames]{xcolor}
\usepackage{listings}
\usepackage[most]{tcolorbox}
\usepackage{fancyhdr}
\usepackage{titlesec}
\usepackage{natbib}
\usepackage{makeidx}
\usepackage{hyperref}

% Fontes distribuídas normalmente com o TeX Live.
\setmainfont{TeX Gyre Pagella}
\setsansfont{TeX Gyre Heros}
\setmonofont{TeX Gyre Cursor}
\setlength{\emergencystretch}{1em}

\definecolor{AzulLivro}{HTML}{183B56}
\definecolor{DouradoLivro}{HTML}{D19A32}
\definecolor{FundoCodigo}{HTML}{F5F7F9}
\definecolor{VerdeCodigo}{HTML}{356859}

\hypersetup{
  colorlinks=true,
  linkcolor=AzulLivro,
  citecolor=VerdeCodigo,
  urlcolor=MidnightBlue,
  pdftitle={LaTeX do Zero ao Primeiro Livro}
}

\pagestyle{fancy}
\fancyhf{}
\fancyhead[LE]{\small\thepage\quad\nouppercase{\leftmark}}
\fancyhead[RO]{\small\nouppercase{\rightmark}\quad\thepage}
\renewcommand{\headrulewidth}{0.4pt}

\titleformat{\chapter}[display]
  {\sffamily\bfseries\color{AzulLivro}}
  {\filleft\Large\chaptername\ \thechapter}
  {1ex}
  {\titlerule\vspace{1ex}\Huge}

\lstset{
  basicstyle=\ttfamily\small,
  backgroundcolor=\color{FundoCodigo},
  keywordstyle=\color{AzulLivro}\bfseries,
  commentstyle=\color{VerdeCodigo}\itshape,
  stringstyle=\color{BrickRed},
  numbers=left,
  numberstyle=\tiny\color{gray},
  frame=single,
  rulecolor=\color{gray!30},
  breaklines=true,
  showstringspaces=false,
  tabsize=2
}

\lstdefinelanguage{LaTeXBook}{
  language=[LaTeX]TeX,
  morekeywords={documentclass,usepackage,begin,end,chapter,section,
    includegraphics,label,ref,cite,frontmatter,mainmatter,backmatter,include}
}

\newtcolorbox{pratica}[1][]{
  colback=blue!3,colframe=AzulLivro,fonttitle=\bfseries,
  title=Pratique agora,#1
}
\newtcolorbox{atencao}[1][]{
  colback=yellow!8,colframe=DouradoLivro!80!black,fonttitle=\bfseries,
  title=Atenção,#1
}
\newtcolorbox{resultado}[1][]{
  colback=green!4,colframe=VerdeCodigo,fonttitle=\bfseries,
  title=Resultado esperado,#1
}

\newcommand{\comando}[1]{\texttt{\textbackslash#1}}
\newcommand{\arquivo}[1]{\texttt{#1}}
\makeindex


\title{LaTeX do Zero ao Primeiro Livro}
\author{Um guia prático e progressivo}
\date{\today}

\begin{document}

\frontmatter
\begin{titlepage}
  \centering
  \vspace*{2.5cm}
  {\sffamily\color{AzulLivro}\Large GUIA PRÁTICO\par}
  \vspace{1cm}
  {\sffamily\bfseries\fontsize{36}{42}\selectfont
    LaTeX do Zero ao\par Primeiro Livro\par}
  \vspace{1cm}
  {\color{DouradoLivro}\rule{0.7\textwidth}{3pt}\par}
  \vspace{1cm}
  {\Large Escreva, organize e publique documentos profissionais\par}
  \vfill
  {\large Edição em desenvolvimento --- \today\par}
\end{titlepage}

\chapter*{Como usar este tutorial}
\addcontentsline{toc}{chapter}{Como usar este tutorial}
Este é um livro para praticar. Digite os exemplos, compile após cada pequena
mudança e compare o resultado. Ao final, você terá entendido a estrutura do
próprio arquivo que está lendo e poderá adaptá-lo para seu primeiro livro.

\tableofcontents

\mainmatter
\chapter{Primeiros passos}
\label{chap:primeiros-passos}

\section{O que é LaTeX?}

LaTeX é um sistema de preparação de documentos. Em vez de posicionar cada
elemento visualmente, você escreve texto e comandos que descrevem sua função:
isto é um título, aquilo é uma seção, esta frase merece ênfase. O mecanismo de
compilação interpreta o arquivo-fonte e produz um PDF consistente.

Essa separação entre conteúdo e apresentação é especialmente útil em livros,
artigos, trabalhos acadêmicos e documentos técnicos. Num projeto longo, mudar a
aparência de todos os capítulos deixa de ser uma tarefa manual: altera-se uma
regra compartilhada. LaTeX\index{LaTeX} também automatiza sumário, numeração, referências
cruzadas, bibliografia e índice remissivo \citep{lamport1994}.

Os arquivos de texto normalmente terminam em \arquivo{.tex}. Eles podem ser
editados no VS Code, TeXstudio, Overleaf ou qualquer editor de texto simples.
Uma distribuição como TeX Live fornece o compilador e os pacotes.

\section{Seu primeiro documento}

Crie uma pasta chamada \arquivo{meu-livro} e, dentro dela, um arquivo
\arquivo{main.tex} com este conteúdo:

\begin{lstlisting}[language=LaTeXBook]
\documentclass{article}

\begin{document}
Ola, LaTeX!
\end{document}
\end{lstlisting}

A declaração \comando{documentclass\{article\}} escolhe a classe do documento.
Ela define decisões básicas de layout e comandos disponíveis. O ambiente
\arquivo{document} delimita o conteúdo que aparecerá no PDF. Tudo antes dele é
o \emph{preâmbulo}, onde configuramos pacotes e opções.

Compile no terminal:

\begin{lstlisting}
latexmk -lualatex main.tex
\end{lstlisting}

O arquivo \arquivo{main.pdf} deve surgir na mesma pasta. O comando
\arquivo{latexmk} executa quantas passagens forem necessárias. A opção
\arquivo{-lualatex} seleciona LuaLaTeX, mecanismo moderno com suporte direto a
Unicode e fontes OpenType.

\begin{atencao}
  Não edite arquivos como \arquivo{.aux}, \arquivo{.log} ou \arquivo{.toc}.
  Eles são gerados durante a compilação. Seu conteúdo pertence aos arquivos
  \arquivo{.tex}.
\end{atencao}

\section{Comandos, argumentos e ambientes}

Um comando começa com barra invertida. Argumentos obrigatórios aparecem entre
chaves e opções aparecem entre colchetes:

\begin{lstlisting}[language=LaTeXBook]
\documentclass[11pt]{book}
\end{lstlisting}

Aqui, \arquivo{book} é obrigatório e \arquivo{11pt} é opcional. Ambientes têm
início e fim explícitos, como \comando{begin\{document\}} e
\comando{end\{document\}}. Se os nomes não coincidirem, a compilação falhará.

Alguns caracteres possuem significado especial: \verb|%| inicia comentário;
\verb|#| representa parâmetros de comandos; \verb|&| separa colunas. Para
imprimi-los no texto, geralmente usamos uma barra: \verb|\%|, \verb|\#| e
\verb|\&|.

\begin{pratica}
  Altere a classe para \arquivo{book}, acrescente a opção \arquivo{12pt} e
  escreva dois parágrafos separados por uma linha vazia. Compile novamente e
  observe o resultado.
\end{pratica}

\begin{resultado}
  Você já consegue identificar preâmbulo, corpo, comando, argumento e ambiente,
  além de transformar um arquivo-fonte em PDF.
\end{resultado}

\chapter{Texto e estrutura}
\label{chap:texto-estrutura}

\section{Parágrafos e ênfase}

No código-fonte, uma quebra de linha comum equivale a um espaço. Uma linha vazia
inicia outro parágrafo. Deixe o LaTeX decidir onde quebrar as linhas; isso mantém
o texto adaptável quando margens ou fontes mudarem.

Use comandos semânticos para destacar conteúdo:

\begin{lstlisting}[language=LaTeXBook]
Este termo e \emph{importante}.
Este trecho esta em \textbf{negrito}.
O comando \LaTeX produz o logotipo do sistema.
\end{lstlisting}

Evite usar negrito e tamanho manual para simular títulos. Os comandos de seção
registram a hierarquia, controlam a numeração e alimentam o sumário automático.

\section{Capítulos, seções e sumário}

A classe \arquivo{book} oferece capítulos. Uma estrutura inicial pode ser:

\begin{lstlisting}[language=LaTeXBook]
\documentclass{book}
\begin{document}
\tableofcontents
\chapter{Introducao}
\section{O problema}
\subsection{Contexto historico}
\chapter{Desenvolvimento}
\end{document}
\end{lstlisting}

O sumário pode aparecer vazio na primeira compilação. LaTeX grava os títulos em
um arquivo auxiliar e os lê na passagem seguinte. \arquivo{latexmk} cuida disso
automaticamente.

Use a hierarquia com intenção. Um capítulo trata de um grande assunto; uma seção
organiza uma etapa desse assunto; uma subseção detalha um aspecto. Muitos níveis
produzem um texto fragmentado e difícil de navegar.

\section{Listas e notas}

Listas não numeradas usam \arquivo{itemize}; sequências usam
\arquivo{enumerate}:

\begin{lstlisting}[language=LaTeXBook]
\begin{itemize}
  \item Planejar o conteudo.
  \item Escrever o primeiro rascunho.
  \item Revisar e compilar.
\end{itemize}

\begin{enumerate}
  \item Criar a pasta.
  \item Criar o arquivo principal.
  \item Gerar o PDF.
\end{enumerate}
\end{lstlisting}

Notas de rodapé são inseridas no próprio ponto da chamada com
\comando{footnote\{Texto da nota.\}}. Use-as para observações complementares,
não para esconder explicações essenciais.

\section{Rótulos e referências cruzadas}

Nunca digite manualmente ``veja o capítulo 3''. A numeração pode mudar. Coloque
um rótulo após o elemento e use \comando{ref}:

\begin{lstlisting}[language=LaTeXBook]
\chapter{Metodos}
\label{chap:metodos}

Como veremos no Capitulo~\ref{chap:metodos}, ...
\end{lstlisting}

O caractere \verb|~| cria um espaço que não permite quebra entre a palavra e o
número. Prefixos como \arquivo{chap:}, \arquivo{sec:}, \arquivo{fig:} e
\arquivo{tab:} ajudam a reconhecer cada tipo de rótulo.

\begin{pratica}
  Crie dois capítulos com duas seções cada. Insira um sumário e faça o segundo
  capítulo referenciar uma seção do primeiro. Depois, inverta a ordem dos
  capítulos e confirme que a referência se atualiza.
\end{pratica}


\chapter{Layout, pacotes e fontes}
\label{chap:layout-fontes}

\section{Estendendo LaTeX com pacotes}

Pacotes adicionam recursos ao documento. Eles são carregados no preâmbulo por
uma declaração própria. O pacote \arquivo{geometry}, por exemplo, configura
papel e margens:

\begin{lstlisting}[language=LaTeXBook]
\usepackage[a4paper,margin=2.5cm]{geometry}
\usepackage{graphicx} % inclusao de imagens
\usepackage{xcolor}  % definicao de cores
\usepackage{hyperref} % links no PDF
\end{lstlisting}

Carregue apenas o necessário e mantenha a configuração centralizada. Em um
livro com vários capítulos, um arquivo \arquivo{preamble.tex} evita que cada
capítulo redefina cores, margens e estilos.

\section{Por que LuaLaTeX?}

LuaLaTeX\index{LuaLaTeX} lê Unicode nativamente e, com \arquivo{fontspec}, usa fontes TrueType e
OpenType do sistema. Isso facilita criar uma identidade visual sem converter
arquivos de fonte para formatos antigos.

\begin{lstlisting}[language=LaTeXBook]
\usepackage{fontspec}
\setmainfont{TeX Gyre Pagella}
\setsansfont{TeX Gyre Heros}
\setmonofont{Latin Modern Mono}
\end{lstlisting}

A fonte principal atende ao texto corrido; a sem serifa é apropriada para
títulos e elementos gráficos; a monoespaçada atende código. Para usar uma fonte
instalada, substitua o nome, por exemplo:

\begin{lstlisting}[language=LaTeXBook]
\setmainfont{Libertinus Serif}
\setsansfont{Source Sans 3}
\setmonofont{JetBrains Mono}
\end{lstlisting}

Confira os nomes reconhecidos pelo sistema com \arquivo{fc-list} no Linux. Uma
mensagem ``font not found'' normalmente indica nome incorreto ou fonte ausente.

\section{Fontes armazenadas no projeto}

Guardar arquivos licenciados para redistribuição numa pasta \arquivo{fonts/}
torna o resultado reproduzível em outras máquinas:

\begin{lstlisting}[language=LaTeXBook]
\setmainfont[
  Path=fonts/,
  UprightFont=MinhaFonte-Regular.otf,
  BoldFont=MinhaFonte-Bold.otf,
  ItalicFont=MinhaFonte-Italic.otf
]{MinhaFonte}
\end{lstlisting}

Verifique sempre a licença. Algumas fontes permitem uso no PDF, mas não a
redistribuição do arquivo original.

\section{Cores e comandos próprios}

Nomes centralizados deixam alterações futuras mais seguras:

\begin{lstlisting}[language=LaTeXBook]
\definecolor{AzulLivro}{HTML}{183B56}
\newcommand{\termo}[1]{\textbf{\color{AzulLivro}#1}}
\end{lstlisting}

Depois, \verb|\termo{tipografia}| aplica sempre o mesmo tratamento. Defina
comandos pelo significado, não por uma aparência passageira. Um nome como
\comando{termo} é mais durável que um nome baseado em cor e tamanho.

\begin{pratica}
  Escolha uma fonte para texto, outra para títulos e uma monoespaçada. Altere
  uma cor e compile. Avalie legibilidade em parágrafos longos, itálico, negrito
  e código antes de confirmar a combinação.
\end{pratica}

\chapter{Imagens, tabelas e matemática}
\label{chap:elementos}

\section{Figuras}

O pacote \arquivo{graphicx} fornece \comando{includegraphics}. Coloque imagens
em uma pasta própria e use o ambiente \arquivo{figure} para legenda, numeração e
referência:

\begin{lstlisting}[language=LaTeXBook]
\begin{figure}[htbp]
  \centering
  \includegraphics[width=0.75\textwidth]{images/diagrama.pdf}
  \caption{Fluxo de producao do livro.}
  \label{fig:fluxo}
\end{figure}
\end{lstlisting}

As letras \arquivo{htbp} indicam posições preferidas: aqui, topo, base ou página
de figuras. São sugestões, pois LaTeX escolhe uma posição que preserve a boa
composição da página. Use PDF para desenhos vetoriais e PNG ou JPEG para imagens
fotográficas. Evite capturas de tela de texto quando o conteúdo puder ser
escrito diretamente.

\section{Tabelas}

O ambiente \arquivo{table} fornece legenda; \arquivo{tabular} organiza células.
O pacote \arquivo{booktabs} cria linhas com proporções profissionais:

\begin{lstlisting}[language=LaTeXBook]
\begin{table}[htbp]
  \centering
  \caption{Classes comuns de documento.}
  \label{tab:classes}
  \begin{tabular}{ll}
    \toprule
    Classe & Uso principal \\
    \midrule
    article & textos curtos \\
    report  & relatorios longos \\
    book    & livros com capitulos \\
    \bottomrule
  \end{tabular}
\end{table}
\end{lstlisting}

Em \verb|{ll}|, cada letra representa uma coluna alinhada à esquerda. Use
\arquivo{c} para centralizar e \arquivo{r} para alinhar à direita. Não use linhas
verticais sem necessidade: espaço e alinhamento geralmente separam melhor os
dados.

\section{Expressões matemáticas}

Matemática curta pode aparecer no parágrafo, como $E=mc^2$. Expressões que
merecem destaque usam um ambiente próprio:

\begin{lstlisting}[language=LaTeXBook]
\begin{equation}
  a^2 + b^2 = c^2
  \label{eq:pitagoras}
\end{equation}
\end{lstlisting}

O pacote \arquivo{amsmath} acrescenta ambientes para alinhar várias expressões:

\begin{lstlisting}[language=LaTeXBook]
\begin{align}
  x + y &= 10 \\
  x - y &= 2
\end{align}
\end{lstlisting}

O símbolo \verb|&| marca o ponto de alinhamento. Não use modo matemático apenas
para obter itálico; isso altera espaçamento e significado dos caracteres.

\begin{pratica}
  Adicione ao seu documento uma imagem com legenda, uma tabela de três linhas e
  uma equação numerada. Crie rótulos e cite os três elementos no texto com
  \comando{ref}.
\end{pratica}


\chapter{Referências e bibliografia}
\label{chap:referencias}

\section{Separando dados bibliográficos}

Referências\index{bibliografia} não devem ser digitadas e numeradas manualmente. Um arquivo
\arquivo{.bib} armazena cada fonte como um registro identificado por uma chave:

\begin{lstlisting}
@book{lamport1994,
  author    = {Leslie Lamport},
  title     = {LaTeX: A Document Preparation System},
  year      = {1994},
  publisher = {Addison-Wesley}
}
\end{lstlisting}

No texto, \comando{citep} produz uma citação entre parênteses. O comando
\comando{citet} incorpora o autor à frase quando usamos \arquivo{natbib}. Ao
final de \arquivo{main.tex}, indicamos o estilo e o banco:

\begin{lstlisting}[language=LaTeXBook]
\bibliographystyle{plainnat}
\bibliography{bibliography}
\end{lstlisting}

Somente fontes citadas aparecem na lista. Use chaves curtas e estáveis, seguindo
o padrão \emph{sobrenome, ano e tema}. Complete autor, título, ano e publicação;
para páginas web, registre endereço e data de acesso.

\section{O ciclo de compilação}

A bibliografia exige mais de uma etapa: LaTeX descobre as citações, BibTeX monta
a lista e LaTeX incorpora e numera o resultado. \arquivo{latexmk} detecta essa
sequência. Se aparecer \verb|[?]|, compile novamente e procure no log por uma
chave inexistente ou erro no arquivo \arquivo{.bib}.

Projetos acadêmicos podem preferir \arquivo{biblatex} com Biber, que oferece
controle mais amplo. Para começar, \arquivo{natbib} e BibTeX formam um fluxo
simples e suficiente. Não misture os dois sistemas na mesma configuração.

\section{Índice remissivo}

O índice remissivo ajuda o leitor a localizar conceitos, não apenas títulos.
Carregue \arquivo{makeidx} e ative a criação do índice no preâmbulo. Marque os
termos no texto e imprima a lista no final do documento.

\begin{lstlisting}[language=LaTeXBook]
LuaLaTeX\index{LuaLaTeX} permite usar fontes OpenType.
\end{lstlisting}

Marque apenas ocorrências úteis. Indexar toda repetição cria ruído. Assim como a
bibliografia, o índice depende de processamento auxiliar que o \arquivo{latexmk}
normalmente executa.

\begin{pratica}
  Adicione um livro e uma página web ao seu \arquivo{bibliography.bib}. Cite-os
  em um parágrafo e marque três conceitos para o índice remissivo. Compile e
  confirme as duas listas no final do PDF.
\end{pratica}

\chapter{Organizando um livro}
\label{chap:organizacao}

\section{Dividir para manter}

Um único arquivo funciona para exemplos curtos, mas dificulta a navegação em um
livro. Mantenha a configuração em \arquivo{preamble.tex}, cada capítulo em
\arquivo{chapters/} e deixe \arquivo{main.tex} mostrar a ordem geral:

\begin{lstlisting}[language=LaTeXBook]
\documentclass{book}
\usepackage{fontspec}
\usepackage[brazil]{babel}
\usepackage{microtype}
\usepackage[a4paper,margin=2.6cm,headheight=15pt]{geometry}
\usepackage{graphicx}
\usepackage{booktabs}
\usepackage{longtable}
\usepackage{amsmath}
\usepackage[dvipsnames]{xcolor}
\usepackage{listings}
\usepackage[most]{tcolorbox}
\usepackage{fancyhdr}
\usepackage{titlesec}
\usepackage{natbib}
\usepackage{makeidx}
\usepackage{hyperref}

% Fontes distribuídas normalmente com o TeX Live.
\setmainfont{TeX Gyre Pagella}
\setsansfont{TeX Gyre Heros}
\setmonofont{TeX Gyre Cursor}
\setlength{\emergencystretch}{1em}

\definecolor{AzulLivro}{HTML}{183B56}
\definecolor{DouradoLivro}{HTML}{D19A32}
\definecolor{FundoCodigo}{HTML}{F5F7F9}
\definecolor{VerdeCodigo}{HTML}{356859}

\hypersetup{
  colorlinks=true,
  linkcolor=AzulLivro,
  citecolor=VerdeCodigo,
  urlcolor=MidnightBlue,
  pdftitle={LaTeX do Zero ao Primeiro Livro}
}

\pagestyle{fancy}
\fancyhf{}
\fancyhead[LE]{\small\thepage\quad\nouppercase{\leftmark}}
\fancyhead[RO]{\small\nouppercase{\rightmark}\quad\thepage}
\renewcommand{\headrulewidth}{0.4pt}

\titleformat{\chapter}[display]
  {\sffamily\bfseries\color{AzulLivro}}
  {\filleft\Large\chaptername\ \thechapter}
  {1ex}
  {\titlerule\vspace{1ex}\Huge}

\lstset{
  basicstyle=\ttfamily\small,
  backgroundcolor=\color{FundoCodigo},
  keywordstyle=\color{AzulLivro}\bfseries,
  commentstyle=\color{VerdeCodigo}\itshape,
  stringstyle=\color{BrickRed},
  numbers=left,
  numberstyle=\tiny\color{gray},
  frame=single,
  rulecolor=\color{gray!30},
  breaklines=true,
  showstringspaces=false,
  tabsize=2
}

\lstdefinelanguage{LaTeXBook}{
  language=[LaTeX]TeX,
  morekeywords={documentclass,usepackage,begin,end,chapter,section,
    includegraphics,label,ref,cite,frontmatter,mainmatter,backmatter,include}
}

\newtcolorbox{pratica}[1][]{
  colback=blue!3,colframe=AzulLivro,fonttitle=\bfseries,
  title=Pratique agora,#1
}
\newtcolorbox{atencao}[1][]{
  colback=yellow!8,colframe=DouradoLivro!80!black,fonttitle=\bfseries,
  title=Atenção,#1
}
\newtcolorbox{resultado}[1][]{
  colback=green!4,colframe=VerdeCodigo,fonttitle=\bfseries,
  title=Resultado esperado,#1
}

\newcommand{\comando}[1]{\texttt{\textbackslash#1}}
\newcommand{\arquivo}[1]{\texttt{#1}}
\makeindex


\begin{document}
\frontmatter
\tableofcontents

\mainmatter
\include{chapters/01-introducao}
\include{chapters/02-fundamentos}

\backmatter
\bibliography{bibliography}
\end{document}
\end{lstlisting}

\comando{input} insere um arquivo no ponto atual. \comando{include} também o
insere, mas inicia uma página e mantém arquivos auxiliares separados, comportamento
conveniente para capítulos. Durante a escrita, \comando{includeonly} compila
apenas partes selecionadas sem perder a numeração conhecida:

\begin{lstlisting}[language=LaTeXBook]
\includeonly{chapters/02-fundamentos}
\end{lstlisting}

\section{Partes de um livro}

A classe \arquivo{book} distingue três regiões. \comando{frontmatter} usa
numeração romana e recebe capa, apresentação e sumário. \comando{mainmatter}
inicia capítulos e numeração arábica. \comando{backmatter} recebe apêndices,
bibliografia e índice.

Uma capa simples pode ser construída com \arquivo{titlepage}. Para páginas sem
numeração visível, use \comando{thispagestyle\{empty\}}. Não tente aperfeiçoar a
capa antes de haver conteúdo suficiente: tipografia e proporções são mais fáceis
de avaliar quando o livro já possui páginas reais.

\section{Convenções de projeto}

Adote nomes previsíveis: \arquivo{01-introducao.tex},
\arquivo{02-fundamentos.tex} e assim por diante. Use letras minúsculas, hífens e
evite espaços. Guarde apenas fontes editáveis no controle de versão; arquivos
\arquivo{.aux}, logs e PDFs gerados entram no \arquivo{.gitignore}.

Cada capítulo deve ter uma finalidade clara. Um modelo simples inclui abertura,
conceitos, exemplo, prática e síntese. A repetição dessa estrutura dá ritmo ao
livro e deixa lacunas evidentes durante a revisão.

\begin{pratica}
  Transforme seu documento em um projeto com \arquivo{main.tex},
  \arquivo{preamble.tex} e dois arquivos em \arquivo{chapters/}. Confirme que o
  PDF continua igual e teste \comando{includeonly}.
\end{pratica}


\chapter{Compilação, revisão e erros}
\label{chap:erros}

\section{Ler a primeira mensagem útil}

Quando uma compilação falha, o terminal pode exibir muito texto. Procure a
primeira linha iniciada por \verb|!| ou uma referência no formato
\arquivo{arquivo.tex:linha}. Erros posteriores frequentemente são apenas
consequências do primeiro.

Problemas comuns incluem:

\begin{itemize}
  \item \emph{Undefined control sequence}: comando digitado incorretamente ou
    pacote ausente;
  \item \emph{File not found}: caminho, extensão ou maiúsculas não coincidem;
  \item \emph{Missing \}}: chave aberta sem fechamento;
  \item \emph{Environment ... undefined}: ambiente incorreto ou pacote não
    carregado;
  \item \emph{There were undefined references}: rótulo ausente ou necessidade
    de outra passagem.
\end{itemize}

Reduza o problema: comente temporariamente partes recentes, compile e reative
pequenos blocos. Não apague arquivos-fonte para ``limpar'' o projeto. Use
\arquivo{latexmk -C} para remover apenas artefatos.

\section{Avisos que merecem revisão}

Uma compilação bem-sucedida ainda pode conter avisos. \emph{Overfull hbox}
indica conteúdo ultrapassando a margem, muitas vezes por URL ou código longo.
Referências e citações indefinidas aparecem como interrogações no PDF. Fontes
substituídas podem alterar a aparência esperada.

Revise também o resultado visual: páginas isoladas, títulos no fim da página,
imagens ilegíveis, inconsistência de termos e espaços excessivos. A ausência de
erros técnicos não garante um bom livro.

\section{Um fluxo diário seguro}

Trabalhe em incrementos pequenos:

\begin{enumerate}
  \item escreva uma seção curta;
  \item compile com \arquivo{latexmk -lualatex main.tex};
  \item corrija o primeiro erro e repita;
  \item confira avisos e páginas alteradas;
  \item registre uma versão funcional no Git.
\end{enumerate}

O modo contínuo, \arquivo{latexmk -lualatex -pvc main.tex}, recompila ao salvar.
No VS Code, a extensão LaTeX Workshop oferece compilação, visualização e salto
entre fonte e PDF.

\begin{atencao}
  Um projeto que depende de LuaLaTeX deve ser compilado sempre com LuaLaTeX.
  Executar pdfLaTeX sobre um preâmbulo com \arquivo{fontspec} produz um erro por
  projeto de forma intencional.
\end{atencao}

\begin{pratica}
  Faça uma cópia de um capítulo e provoque três erros: remova uma chave, altere
  o nome de um ambiente e referencie uma imagem inexistente. Identifique a
  primeira mensagem útil de cada caso e depois restaure o arquivo.
\end{pratica}


\chapter{Projeto final: seu primeiro livro}
\label{chap:projeto-final}

Chegou a hora de transformar as técnicas anteriores num projeto autoral. Use o
modelo da pasta \arquivo{examples/} como ponto de partida ou copie este próprio
projeto e remova o conteúdo didático.

\section{Etapa 1: definir o livro}

Antes do LaTeX, escreva uma frase para cada item:

\begin{itemize}
  \item tema e transformação prometida ao leitor;
  \item público e conhecimento prévio necessário;
  \item entre cinco e dez capítulos, cada um com um resultado observável;
  \item exemplos, imagens ou fontes de pesquisa necessários.
\end{itemize}

Um sumário provisório é uma hipótese, não um contrato. Ele pode mudar quando a
escrita revelar uma ordem mais clara.

\section{Etapa 2: montar a estrutura}

Crie \arquivo{main.tex}, \arquivo{preamble.tex}, \arquivo{bibliography.bib} e as
pastas \arquivo{chapters/} e \arquivo{images/}. Faça um arquivo por capítulo e
adicione todos ao principal com \comando{include}. Compile enquanto os capítulos
ainda contêm apenas títulos. Assim, problemas de infraestrutura são resolvidos
antes de se misturarem a centenas de páginas.

\section{Etapa 3: escolher a identidade visual}

Defina uma fonte de leitura, uma fonte de títulos e uma fonte monoespaçada se
houver código. Escolha uma cor principal e uma cor de destaque. Teste uma página
com título, parágrafos, lista, negrito e itálico. Priorize legibilidade; efeitos
visuais devem apoiar a hierarquia, não competir com o conteúdo.

\section{Etapa 4: escrever e revisar}

Produza primeiro uma versão completa, mesmo imperfeita. Depois faça revisões
separadas:

\begin{enumerate}
  \item \textbf{estrutura}: ordem e objetivo de cada capítulo;
  \item \textbf{conteúdo}: precisão, exemplos e lacunas;
  \item \textbf{linguagem}: clareza, repetição e terminologia;
  \item \textbf{visual}: quebras, figuras, tabelas e consistência;
  \item \textbf{técnica}: log limpo, links, citações e índice.
\end{enumerate}

\section{Lista de verificação}

Seu primeiro livro está pronto para uma versão de leitura quando:

\begin{itemize}
  \item compila do zero com um único comando;
  \item não possui referências ou citações indefinidas;
  \item explica a quem se destina e como deve ser usado;
  \item mantém títulos, exemplos e termos consistentes;
  \item possui fontes e imagens com licenças adequadas;
  \item foi lido no PDF, não apenas no código-fonte;
  \item foi testado por pelo menos uma pessoa do público pretendido.
\end{itemize}

\begin{resultado}
  Ao concluir essas etapas, você terá um livro modular, compilável e preparado
  para crescer. Mais importante: saberá localizar cada decisão editorial no
  código e poderá alterá-la sem refazer o documento inteiro.
\end{resultado}

\section*{Próximos passos}

Explore a documentação oficial do projeto LaTeX \citep{latex-project} à medida
que surgirem necessidades concretas. Aprender pacotes sob demanda é mais
eficiente do que tentar memorizar o ecossistema. Seu primeiro livro não precisa
usar todos os recursos; precisa servir bem ao leitor.


\backmatter
\bibliographystyle{plainnat}
\bibliography{bibliography}
\printindex

\end{document}
